\documentclass[10pt,a4paper]{amsart}
\usepackage[margin=19mm]{geometry}
\usepackage{amsmath,amssymb,mathtools}
\usepackage[scr=rsfs,cal=euler]{mathalfa}
\usepackage{xfrac}
\usepackage[unicode,hidelinks,bookmarks=false]{hyperref}
\newcommand{\ie}{i.e.\ }
\newcommand{\cf}{cf.\ }
\newcommand{\resp}{respectively\ }
\newcommand{\Subsection}{\subsection}
\providecommand{\nobreakdash}{\nobreak\hbox{-}}
\setlength{\emergencystretch}{2em}
\multlinegap0pt
\usepackage{mathtools}
\usepackage{amssymb}
\usepackage{xfrac}
\usepackage{xcolor}
\usepackage{mdframed}
\usepackage{float}
\usepackage[shortlabels]{enumitem}
\usepackage{tikz}
\usepackage{bm}
\newtheorem{proposition}{Proposition}
\newtheorem{lemma}{Lemma}
\newcommand\bbZ{\ensuremath{\mathbb{Z}}}
\newcommand{\ora}[1]{\overrightarrow{#1}}
\newcommand{\reducesto}{\mathrel{\searrow}}
\newcommand{\wid}{\operatorname{wid}}
\title{The intersection dual of geodesic currents}
\author{D\'idac Mart\'inez-Granado, Dylan P. Thurston}
\date{Source: arXiv:2605.04031v1 (CC BY 4.0)}
\begin{document}
\maketitle

\noindent\emph{Source and licence.} The intersection dual of geodesic currents. Version arXiv:2605.04031v1 (CC BY 4.0), distributed by arXiv under the Creative Commons Attribution 4.0 International licence, http://creativecommons.org/licenses/by/4.0/, as recorded in the arXiv licence metadata for this exact version and shown on the abstract page https://arxiv.org/abs/2605.04031v1. Full licence notice: \texttt{licenses/arxiv-2605.04031-CC-BY-4.0.txt}.\par\smallskip

\noindent\textit{Editorial scope.} Counted unit: the article's proposition prop:join-split-geom (source0.tex lines 2965--2987). The article's complete original proof of it is delivered with it (source0.tex lines 3043--3118, one \texttt{proof} environment of 76 lines, which the article places away from the statement). No mathematics is written, repaired or re-derived: the statement and the proof are the article's own lines, in the article's own notation, and no retained line is substituted or re-flowed.  Retained uncounted as the source-local context the proof needs: lines 3007--3016.  Nothing was omitted from any retained span, so the package carries no omission record.  No retained line contains a citation, so the wrapper carries no bibliography and no bibliography entry.  No retained line contains a figure environment, an image-inclusion command or drawing code, so no image is delivered and the package declares no asset dependency.  Editorial formatting only, in addition: an AMS \texttt{amsart} preamble that restates the article's own declarations and macros used by the retained lines, namely the environments \texttt{proposition}, \texttt{lemma} on separate counters, together with the standard packages those lines need.  The retained environments print in this document as Proposition 1, Lemma 1 in order of appearance, whereas the article prints the same items with the numbering of its own full document.  The complete native source of the article as distributed in the arXiv e-print for this exact version is bundled unchanged as \texttt{sources/199832/source0.tex}.
\begin{lemma}
  Given $a,b, x \in \pi_1(S)$ so that $\ora{a}$ and $\ora{b}$ cross $\ora{x}$ to
  the right and $g_1 = x^k ax^\ell$ with $k,\ell \in \bbZ$, we have
  $g_1^+ \in w(x^k a\cdot\ora{x})$ and
  $g_1^- \in w(x^{-\ell}A\cdot\ora{x})$. Similarly for
  $g_2 = x^k a x^\ell b x^n$ with $k,\ell,n\in\bbZ$, we have
  $g_2^+ \in w(x^ka \cdot \ora{x})$ and $g_2^- \in w(x^{-n} B
  \cdot\ora{x})$.
\label{lem:windows}
\end{lemma}

\begin{proposition}\label{prop:join-split-geom}
Let $a,b,x \in \pi_1(S)$ be curves with $\ora{a}$, $\ora{b}$
crossing $\ora{x}$ to the right. Then there is $s \ge 0$ (depending on $a,b,x$)
  so that, for all
 $n, m, r \in \mathbb{Z}$ with either $s < r < m-n-s$ or $m-n+s < r <
 -s$, we have
  \begin{align}
    \label{eq:cross-join} [a x^n] [b x^m] &\reducesto [ax^{n+r}bx^{m-r}]\\
    \label{eq:cross-split} [a x^nb x^m] &\reducesto [ax^{n+r}][bx^{m-r}].
  \end{align}

  There exists $s' \ge 0$ (depending on $a,x$) so that for all
  $|n-m| > s'$,
  \begin{align}
    \label{eq:cross-cancel}[ax^n][ax^m] &\reducesto [x^{|n-m|}].
  \end{align}

 If $x$ is simple and $a=b$, we may take $s=s'=0$. In particular, for
 all $n\in\bbZ$,
 \begin{align}
    [a x^{n+1}] [a x^{n-1}] &\reducesto 2[ax^{n}].\label{eq:cross-simple}
  \end{align}
\end{proposition}

\begin{proof}[Proof of Proposition~\ref{prop:join-split-geom}]
  In each case, we will show that we can choose lifts of the curves to
  the normalized band model for which a
  relevant crossing is essential. We start with
  Eq.~\eqref{eq:cross-join}, looking first at the windows $w(a\cdot
  \ora{x})$ and $w(x^rb\cdot\ora{x})$.
  Since each window is a
  finite size, there is a finite range $[s_1,s_2]$ so that $w(x^rb\cdot\ora{x})$
  is to the left of (resp.~to the right of) $w(a\cdot\ora{x})$ if $r < s_1$
  (resp.\ $r > s_2$).
  Similarly there
  is a range $[s_1',s_2']$ so that $w(A\cdot\ora{x})$ is disjoint from
  $w(x^rB\cdot\ora{x})$ if $r \notin [s_1',s_2']$, and the exponents determine the
  ordering of the windows. Set
  $s \coloneqq \max\{s_2,s_2',-s_1,-s_1'\}$.

  Now consider
  $g_1 = ax^n$ and $g_2 = x^rbx^{m-r}$, with $s < r < m-n-s$. We have
  \[
    r-m < -n-s \le -n+s_1',
  \]
  so, by Lemma~\ref{lem:windows}, $g_1^- \in w(x^{r-m}B\cdot\ora{x})$
  is to the left of $g_2^- \in w(x^{-n}A\cdot\ora{x})$, and
  \[
    s_2 \le s < r,
  \]
  so $g_1^+ \in w(x^rb\cdot\ora{x})$ is to the right of $g_2^+ \in
  w(a\cdot\ora{x})$. Thus
  $\ora{x^rbx^{m-r}}$ crosses $\ora{ax^n}$ to the left. Similarly if
  $m-n+s < r < -s$, then $\ora{x^rbx^{m-r}}$ crosses $\ora{ax^n}$ to the
  right. In either case,
  \[
    [ax^n][bx^m] = [ax^n][x^r bx^{m-r}] \reducesto [ax^{n+r}bx^{m-r}]
  \]
  by oriented disconnected smoothing.
  For Equation~\eqref{eq:cross-split}, write $[x^r ax^n bx^{m-r}]=[gh]$ with
  $g=x^r ax^n$ and $h=bx^{m-r}$, 
  so $g^+ \in w(x^ra\cdot\ora{x})$ and $g^- \in w(x^{-n}A \cdot \ora{x})$, and $h^+ \in w(b \cdot \ora{x})$, $h^- \in w(x^{r-m} B \cdot \ora{x})$.
  Hence, if $s < r < m-n-s$, we have $s_1 < r$
  so $w(b\cdot\ora{x})$ is to the left of $w(x^r a\cdot\ora{x})$,
  while $r-m < -n+s_1'$, so $w(x^{r-m}B \cdot\ora{x})$ is to the left
  of $w(x^{-n}A\cdot\ora{x})$, and hence $\ora{h}$ and $\ora{g}$ are R-parallel, giving Eq.~\eqref{eq:cross-split} by oriented connected smoothing.
  
  We similarly get parallelism if $m-n+s < r < -s$.

  To see Eq.~\eqref{eq:cross-cancel}, suppose without loss of generality
  that $n < m$. Let $s_0 = \lceil \wid(w(a\cdot\ora{x})) \rceil$
  (with width measured in the normalized band model),
  and
  let $s' = 2s_0 - 1$. Since $m-n \ge 2s_0$, we can pick $r$ so that
  $s_0 \le r \le m-n-s_0$. By the above argument, $\ora{ax^n}$ and
  $\ora{x^rax^{m-r}}$ R-cross (since, e.g., $(ax^n)^+ \in w(a\cdot\ora{x})$ and
  $(x^rax^{m-r})^+ \in w(x^r a\cdot\ora{x})$, which are disjoint), so we get the
  desired smoothing in Eq.~\eqref{eq:cross-cancel} by using unoriented
  disconnected smoothing.

  In the cases when $x$ is simple, two windows are disjoint, nested,
  or identical; this implies that $s_1 = s_2$ and $s_1' = s_2'$.
  If in
  addition $a=b$, we also see that $s_1=s_1'=0=s$. This gives Eqs.\
  \eqref{eq:cross-join} and \eqref{eq:cross-split}, and the special
  case Eq.\ \eqref{eq:cross-simple}.

  To see that with $x$ simple we get Eq.\ \eqref{eq:cross-cancel} with
  $s'=0$ (i.e., $[ax^n][ax^{n+1}] \reducesto [x]$), we look slightly deeper
  in the words. The analysis above shows $(ax^n)^- \in
  w(x^{-n}A\cdot\ora{x})$, which is to the right of
  $(ax^{n+1})^- \in w(x^{-n-1}A\cdot\ora{x})$. By
  looking at the band model centered on $a \cdot \ora{x}$ (or
  equivalently by conjugating by $a$), we see that
  $(ax^n)^+ \in w(ax^na\cdot\ora{x}) \subset w(a\cdot\ora{x})$. Similarly
  $(ax^{n+1})^+ \in w(ax^{n+1}a\cdot\ora{x}) \subset
  w(a\cdot\ora{x})$. We therefore have that $(ax^n)^+$ is to the left
  of $(ax^{n+1})^+$, implying that $\ora{ax^n}$ crosses $\ora{ax^{n+1}}$ as
  desired.
\end{proof}

\end{document}
